\begin{annexitem}{annex:p7.eqlk.alpha}{thm:p07-eqlk-alpha}{Literal alpha-conversion kernels with identity atoms}
\begin{proof}
Terms have no bound variables. Alpha-equivalent atomic formulas are
therefore literally identical, including identity atoms; use Id. Falsum
uses Id as well. We prove both directions together by induction on formula
complexity. For \(\neg A\equiv_\alpha\neg B\), start with the induction
kernel \(B\Rightarrow A\), apply \(\neg L\) to its succedent \(A\),
then \(\neg R\) to \(B\), obtaining
\(\neg A\Rightarrow\neg B\). Interchange the letters for the other direction.

For conjunction, weaken the two induction kernels to have both source
conjuncts on the left, apply \(\land L\) in each branch, then
\(\land R\) to obtain the target conjunction. For disjunction, inject
each target disjunct by right weakening and \(\lor R\), then use
\(\lor L\) on the two source-disjunct branches. These shared-context
rules use exactly one source compound on the left of each premise.

For \(A_1\to A_2\equiv_\alpha B_1\to B_2\), the induction kernels
are \(B_1\Rightarrow A_1\) and \(A_2\Rightarrow B_2\).
Right weakening of the first and left weakening of the second give
\(B_1\Rightarrow B_2,A_1\) and
\(B_1,A_2\Rightarrow B_2\). The shared-context \(\to L\) yields
\(B_1,A_1\to A_2\Rightarrow B_2\); \(\to R\) yields
\(A_1\to A_2\Rightarrow B_1\to B_2\). The reverse uses the
opposite induction kernels.

At a quantifier, write the two formulas as \(QxA\) and \(QyB\).
Choose a variable \(z\) occurring nowhere in either formula. Their opened
matrices \(A[z/x]\) and \(B[z/y]\) are alpha-equivalent. To justify
this syntactic step, inspect the generating alpha-renaming relation of
\cref{def:p02-fol-alpha-equivalence}: opening sends the occurrences governed
by the displayed outer binder to the same fresh variable; an inner
renaming is unchanged outside its governed occurrences, and congruence,
symmetry and transitivity preserve this property. Equivalently, assign
each bound occurrence its binder's depth and keep free variable names;
the alpha-generating steps preserve that finite tree, and opening changes
exactly its outer-bound occurrences to free \(z\). The converse comparison
is constructive: list the finitely many binder positions in traversal order,
choose distinct new variable names absent from both original formulas and
their free variables, and rename the binders to those names successively.
Each renaming is a generating alpha step, since its name is absent from
the current matrix. Two formulas with the same binding tree become the
same literal formula by this procedure. Apply it to the two opened matrices
to obtain their alpha-equivalence. Substitution preserves complexity by
\cref{lem:substitution-complexity}, so the matrices fall under the induction.

For \(Q=\forall\), the induction kernel is
\(A[z/x]\Rightarrow B[z/y]\). Apply \(\forall L\) to the source
instance and then \(\forall R\) to the target. The eigenvariable
\(z\) is absent from the source quantified formula and the target
conclusion. For \(Q=\exists\), apply \(\exists R\) to the target
instance first and then \(\exists L\) to the source. The same absence
condition verifies its eigenvariable proviso. Interchanging source and
target gives the reverse kernel. No step used Cut, Mix or an alpha rule.

Finally, to replace a succedent occurrence of \(A\) by \(B\), cut
the existing proof with \(A\Rightarrow B\). To replace an antecedent
occurrence, cut \(B\Rightarrow A\) with that proof. The other contexts
are unchanged and the selected occurrence alone is cut. This last
conversion procedure is deliberately cut-admitting.
\end{proof}
\end{annexitem}

\begin{annexitem}{annex:p7.eqlk.replacement}{thm:p07-eqlk-replacement}{The exact oriented equality kernels}
\begin{proof}
Choose a variable \(z\) absent from \(r,t\). Apply the registered left
replacement to the identity \(t=r\Rightarrow t=r\), using the auxiliary
formula \(t=z\) and the equality \(r=t\). The result is
\[
                       r=t,t=t\Rightarrow t=r.
\]
Cut the equality initial \(\Rightarrow t=t\) against its displayed
\(t=t\) antecedent occurrence. This yields \(r=t\Rightarrow t=r\).
Both Cut premises are Cut/Mix-free. Freshness makes the auxiliary
instances literally \(t=r\) and \(t=t\); no capture or hidden term
substitution is involved.

For arbitrary \(A,x\), apply the primitive equality rule directly to
\(A[r/x]\Rightarrow A[r/x]\). It gives
\(r=t,A[t/x]\Rightarrow A[r/x]\), the middle kernel.
Use the same construction with the roles of \(r,t\) exchanged to get
\(t=r,A[r/x]\Rightarrow A[t/x]\). Cut the symmetry kernel against
its \(t=r\) occurrence to obtain
\(r=t,A[r/x]\Rightarrow A[t/x]\).

If a premise derivation instead ends in
\(\Gamma,A[t/x]\Rightarrow\Delta\), the primitive rule with exchanged
terms gives \(\Gamma,t=r,A[r/x]\Rightarrow\Delta\).
Cut with symmetry to obtain
\(\Gamma,r=t,A[r/x]\Rightarrow\Delta\), the opposite-oriented
left replacement. These operations do not alter \(\Gamma\) or
\(\Delta\). Zero occurrences of \(x\) merely make both instances
truth-equivalent, although the deterministic substitution can freshen an
otherwise irrelevant bound name differently in the two instances. The
registered rule uses those literal instances, so no conversion is presumed.
Multiple occurrences are replaced simultaneously, exactly as the sheet
permits. Occurrences of \(x\) inside \(r\) or \(t\) do not
trigger a second substitution: the source operation inserts a term once.
\end{proof}
\end{annexitem}

\begin{annexitem}{annex:p7.eqlk.congruence}{thm:p07-eqlk-congruence}{Finite explicit identity congruence constructions}
\begin{proof}
Choose a fresh hole variable \(z\). The right-directed replacement kernel
for equality \(t=u\) and formula \(r=z\) is
\(t=u,r=t\Rightarrow r=u\), proving transitivity by multiset exchange.

For the function case start with reflexivity of
\(f(r_1,\ldots,r_n)\). At stage \(j\), use the replacement kernel
for \(r_j=t_j\) with the formula
\[
 f(r_1,\ldots,r_n)
   =f(t_1,\ldots,t_{j-1},z,r_{j+1},\ldots,r_n).
\]
The hole variable is absent from all displayed terms. Its two instances
are exactly the stage \(j-1\) and stage \(j\) equalities. Cut the current
proof against the first instance in the kernel. The conclusion retains
one copy of each equality assumption used through stage \(j\), and
its succedent is the next stage equality. Iterate through the finite
arity; left contraction removes repeated assumptions if two equalities
coincide literally. Weakening restores any required repeated occurrence
in the stated endsequent. For zero arity there is no stage: the initial
is \(\Rightarrow f=f\).

For a relation, start with \(R(r_1,\ldots,r_n)\Rightarrow
R(r_1,\ldots,r_n)\). At stage \(j\) use the formula
\(R(t_1,\ldots,t_{j-1},z,r_{j+1},\ldots,r_n)\), and compose
its replacement kernel with the current proof by Cut. This preserves
the one original relation premise and produces the replaced argument
tuple after \(n\) stages. For the reverse direction use the same
stage construction on the reversed tuples and cut the individual
symmetry kernels for \(r_j=t_j\) into the equality premises it requires.
Contract or weaken only after specifying the required multiset
occurrences. For \(n=0\), both directions are identity of \(R\).
All terms are finite and each stage is a previously constructed finite
kernel; hence these descriptions give terminating proof constructions,
not appeals to semantic substitutivity or to the desired completeness.
\end{proof}
\end{annexitem}

\begin{annexitem}{annex:p7.eqlk.soundness}{thm:p07-eqlk-soundness}{Every structural, logical and equality soundness case}
\begin{proof}
Use induction on the finite derivation tree, proving validity under every
structure and assignment. A counterevaluation of a sequent makes every
antecedent true and every succedent false. The contexted identity leaf
has a formula both true and false at a counterevaluation; the falsum
leaf would make falsum true. Reflexivity holds because each total term
has a value and that value is identical to itself.

For either weakening, a counterevaluation of the conclusion is already
one of the premise. A duplicated formula and a single copy have the
same truth condition, proving both contractions. For Cut, at a
counterevaluation of its conclusion the shared cut formula must be true
by the left premise and false by the right premise. For Mix, any positive
number of copies has the same truth condition as one; the same argument
uses the selected copies while every unselected context occurrence stays
in its declared side. Thus arbitrary multiplicities and empty retained
contexts are covered.

For \(\neg L\), a true principal negation makes its matrix false,
falsifying the premise. For \(\neg R\), a false principal negation
makes its matrix true, again falsifying the premise. For \(\land L\),
truth of the conjunction makes both antecedent matrices true. For
\(\land R\), falsity of the conjunction makes at least one right
matrix false, giving a counterevaluation of the corresponding premise.
For \(\lor L\), truth of the disjunction makes one antecedent matrix
true, giving a counterevaluation of that premise. For \(\lor R\),
falsity makes both right matrices false in its sole premise.
For \(\to L\), truth of the antecedent implication either makes its
antecedent matrix false, falsifying the first premise, or makes its
consequent matrix true, falsifying the second. The contexts are the
same \(\Gamma,\Delta\) in both branches, exactly as on this sheet.
For \(\to R\), falsity of the principal implication makes its matrix
antecedent true and consequent false, falsifying the premise.

The term-substitution lemma \ref{lem:term-substitution} treats variables,
constants and every finite-arity function application. The formula
substitution theorem \ref{thm:fol-substitution} treats both predicate and
identity atoms, falsum and every connective; at a quantifier it treats
the bound-variable-equals-substitution-variable case, the commuting
case and the fresh-renaming case. Its totality therefore supplies the
actual instances below even if a term contains the variable replaced.
There is no additional unproved free-for proviso.

For \(\forall L\), a true quantified antecedent is true at the value
of the supplied term, so the substitution theorem makes its premise
instance true. For \(\exists R\), a false quantified succedent is
false at every value, including that term value. In either case a
counterevaluation of the conclusion therefore falsifies the premise.

For \(\forall R\), a counterevaluation of the conclusion supplies an
element \(d\) falsifying the quantified matrix. Replace the eigenvariable
\(y\)'s assignment by \(d\). It is free in no context formula or
quantified conclusion, so \cref{lem:fol-coincidence} preserves their truth
values. The substitution theorem makes the displayed \(A[y/x]\)
premise false. This includes \(y=x\): its quantified assignment simply
overrides the original value of \(x\).
For \(\exists L\), choose an element witnessing its true quantified
antecedent and perform the same eigenvariable update. Context truth and
all succedent falsity persist by its freshness proviso, while the
substitution theorem makes the displayed antecedent instance true.
Both arguments concern one arbitrary element, not a choice function.

Finally, at a counterevaluation of an \(=L\) conclusion, \(r=t\)
is true, so \(r,t\) have equal values. The substitution theorem evaluates
\(A[r/x]\) and \(A[t/x]\) by the same update of the original
assignment. Thus truth of the conclusion's \(A[t/x]\) makes the
premise's \(A[r/x]\) true, without changing either context. This covers
zero or multiple free occurrences, quantified matrices and arbitrary
terms. Every rule on the existing sheet has now been treated.
\end{proof}
\end{annexitem}

\begin{annexitem}{annex:p7.eqlk.nd}{thm:p07-eqlk-nd}{Complete two-way proof-tree constructions}
\begin{proof}
We first specify the finite context operations used in the constructions.
Natural deduction forgets multiset multiplicity, so every sequent antecedent
becomes its finite support. A disjunction is built with the fixed association
of \(D_\Delta\). From a selected disjunct, repeated \(\lor I\)
injects it at its place in this tree; nested \(\lor E\) eliminates the
tree one leaf at a time. If the support is empty, its formula is \(\bot\)
and \(\bot E\) supplies every target. These operations also reassociate
disjunction trees, delete repeated disjuncts and inject a smaller support
into a larger one by a finite derivation. When a rule description writes
\(D\lor A\), these operations convert its actual finite-support
disjunction to that displayed form and back. In the empty retained-support
case \(A\) injects into \(\bot\lor A\), whose falsum branch is
eliminated by \(\bot E\); repeated or already retained \(A\)'s use
the same injection and branch elimination. They are not literal identifications.

To keep an assumption \(A\) genuinely present when a branch derives
\(C\) without using it, form \(A\land C\) from an actual \(A\)
leaf and that derivation, then project \(C\). Mark the \(A\) leaf
with the branch's discharge label. All new labels are chosen fresh.
This finite detour makes every intended \(\to I,\neg I,\lor E\)
or \(\exists E\) discharge nonvacuous without changing its conclusion.
Grafting a supplied derivation into an open assumption leaf preserves
every old discharged class, because such a leaf is not governed by an
old discharge; freshen inserted labels to avoid collisions.

We also need a literal variable-renaming construction for already obtained
ND proofs. Let \(\pi\) be a permutation of the individual variable
alphabet. Renaming every bound and free variable by \(\pi\) commutes
literally with term formation and each propositional constructor. It
commutes with total formula substitution up to alpha-equivalence:
\[
 \pi(A[t/x])\equiv_\alpha (\pi A)[\pi t/\pi x].
\]
Here is the full substitution check. The variable and function-term clauses
commute literally, hence so do predicate and identity atoms. Falsum and
all connectives inherit the property recursively. At a quantifier whose
binder is the variable replaced, both substitutions stop. If the binder
is not in the inserted term, both traverse the matrix. In the remaining
case both rename the binder to an absent name before recursing; \(\pi\)
preserves absence but need not preserve the order defining the least fresh
name. Their outputs have the same binding tree and free variables, so
the explicit canonical freshening above proves alpha-equivalence.

Induct on the ND proof tree to realize its renamed proof. Assumptions,
reflexivity, propositional rules and discharge labels are renamed literally.
Quantifier or equality instances are first converted by the primitive ND
alpha rule to the required renamed total instance, the same source rule
is applied, and its result is converted to the renamed original conclusion
when necessary. Universal introduction, existential elimination and their
all-occurrence or free-occurrence freshness provisos are preserved by the
bijection; alpha conversion has the same free variables and introduces
no new open assumption. Thus every ND rule, including the explicit alpha
rule, is covered. The proof is finite, and every original marked discharge
leaf remains present, with its formula renamed consistently.

At a strong LK quantifier rule choose a variable \(z\) absent from the
already translated finite ND premise proof and the original LK conclusion,
and transpose its eigenvariable
\(y\) with \(z\) in the already translated ND premise. This is not a
recursive call on a larger LK proof. Since \(y\) is free in no surviving
context formula, each renamed surviving assumption is alpha-equivalent
to its original: replace its open leaf by the original assumption followed
by ND alpha conversion. Convert the renamed succedent tree back in the
same way. The distinguished matrix becomes literally \(A[z/x]\)
after alpha conversion by the substitution check. Consequently the
new \(z\) is absent altogether, not only free-absent, from the original
surviving assumptions, original succedent and original quantified formula.
This supplies the stronger ND existential proviso even when \(y\)
originally occurred bound in an LK context. If an eigenvariable assumption
would be unused, the preceding detour supplies an actual occurrence.

Induct on an equality-LK derivation to build a natural-deduction proof of
\(D_\Delta\) from \(\operatorname{supp}(\Gamma)\).
Id injects its displayed antecedent formula into the succedent disjunction;
\(\bot L\) eliminates its falsum assumption. Left weakening adds an
allowed assumption (use the detour if its actual presence is needed later);
right weakening is disjunction injection. Left and right contraction
alter only multiplicity, not the finite support. For Cut, eliminate the
translated left-premise disjunction: each retained right formula injects
into the desired conclusion, and the cut-formula branch uses the translated
right premise and injects its conclusion. Any unused discharged branch
assumption receives the detour. Mix uses that same translated branch
construction: its selected copies form one formula in the finite support,
while unselected equal copies, if present, remain in the prescribed
retained supports. The finite disjunction conversions above expose the
chosen branch even when its formula also occurs in the retained succedent.
No induction call on a newly constructed Cut proof is needed.

For \(\land L\), project both conjuncts from the compound assumption
and graft them into the translated premise. For \(\lor L\), eliminate
the compound assumption using its two translated branches. For \(\to L\),
eliminate the translated first premise: a retained succedent branch is
already a desired answer, and the antecedent-matrix branch uses the
implication assumption to obtain the consequent matrix, then applies
the translated second premise. For \(\neg L\), the matrix branch
contradicts the negation assumption and uses \(\bot E\).
For \(\forall L\), instantiate the universal assumption by \(\forall E\)
with the exact total instance and graft it into the premise.
For \(\exists L\), first freshen the already translated premise by
the preceding ND construction. Eliminate the original existential
assumption with that \(A[z/x]\) minor branch. Its new eigenvariable
is absent altogether from the surviving assumptions, conclusion and
existential formula, as the exact ND sheet requires. Add the branch
detour if needed.

For \(\lor R\), each of the two principal disjunct branches injects
into the new compound; retained branches inject unchanged. For
\(\exists R\), the witness-instance branch applies \(\exists I\),
and retained branches again inject unchanged. For \(\land R\), eliminate
the translated first premise: a retained branch is finished, while
its first-matrix branch eliminates the translated second premise;
combine its second-matrix branch by \(\land I\), leaving retained
branches unchanged. This is a finite pair of case eliminations and
uses both common-context premises exactly as registered.

Three remaining right cases use classical RAA. For \(\neg R\),
write \(D\) for the retained succedent disjunction. The premise
translation proves \(D\) under \(A\). Assume
\(\neg(D\lor\neg A)\). Injection of an assumed \(D\) into the
forbidden disjunction yields \(\neg D\); injection of an assumed
\(\neg A\) yields \(\neg\neg A\). RAA gives \(A\), the premise
proof gives \(D\), and the contradiction closes the outer RAA,
proving \(D\lor\neg A\). Finally use the disjunction-support operations
to reach the specified \(D_\Delta\).
For \(\to R\), assume \(\neg(D\lor(A\to B))\) and derive
\(\neg D\) by the same injection. Under an actual \(A\) assumption,
eliminate the translated premise \(D\lor B\): its \(D\)-branch
contradicts \(\neg D\) and yields \(B\); its other branch gives \(B\).
Discharge \(A\), using the detour if needed, inject \(A\to B\)
into the forbidden disjunction, and apply outer RAA.
For \(\forall R\), first freshen its translated premise to
\(D\lor A[z/x]\) by the preceding construction. Assume
\(\neg(D\lor\forall xA)\) and derive
\(\neg D\). Eliminate the translated disjunction to obtain
\(A[z/x]\); the \(D\)-branch uses contradiction and \(\bot E\).
The eigenvariable \(z\) is absent from every surviving assumption,
including the temporary outer negation. Apply \(\forall I\) to \(z\), then
the explicit ND alpha rule to obtain literally \(\forall xA\),
inject it and close RAA. The two quantified formulas are alpha-equivalent:
\(z\) is absent from the original quantified formula, and a fresh
common opening verifies their matrices, as in Annex~\ref{annex:p7.eqlk.alpha}.
These cases work with \(D=\bot\) as well. Every negation or RAA
assumption is used in the displayed contradiction.

Equality introduces two cases, not a new semantic shortcut. The equality
initial translates by ND \(=I\). For \(=L\), the translated premise
requires \(A[r/x]\). From the new assumption \(r=t\), ND symmetry
gives \(t=r\): start with reflexivity \(r=r\), and use ND replacement
in the fresh-hole formula \(z=r\). Then \(t=r\) and the new assumption
\(A[t/x]\) give \(A[r/x]\) by ND \(=E\). Graft this derivation
into every required open \(A[r/x]\) leaf of the premise proof. The
total instances are exactly those on the sheet, including terms with
the hole variable; they insert once, rather than recursively substituting
inside the inserted term. This completes every LK rule case.

For the reverse direction, induct on a derivation on the exact ND sheet.
At each node use one antecedent occurrence for each formula in its
surviving open-assumption support. To align two or three premises,
weaken their contexts to the union of these supports. Temporary
discharge formulas have one additional selected antecedent occurrence,
even if an unmarked occurrence of the same formula survives. After
Cut, contract duplicate surviving occurrences to one; at the end
weakening restores any requested multiplicities.

An assumption translates to contexted Id. Conjunction introduction is
\(\land R\) after context alignment; conjunction elimination cuts
the premise against \(A\land B\Rightarrow A\) or
\(A\land B\Rightarrow B\), each obtained by \(\land L\) over Id.
Disjunction introduction weakens by the missing disjunct and uses
\(\lor R\). For disjunction elimination, align the minor contexts,
apply \(\lor L\) to their selected temporary assumptions, then cut
the translated major premise into the resulting \(A\lor B\) occurrence.
Exactly the assumptions governed by the source labels are removed;
unmarked equal formulas survive in the aligned context.

Implication introduction is \(\to R\) on the selected assumption.
For implication elimination derive \(A,A\to B\Rightarrow B\):
its shared \(\to L\) premises are \(A\Rightarrow B,A\) and
\(A,B\Rightarrow B\), both contexted identities or weakenings of Id.
Cut the two supplied proofs into \(A\to B\) and \(A\).
Negation introduction first cuts a falsum succedent against
\(\bot\Rightarrow\), then applies \(\neg R\) to the selected
assumption. Negation elimination cuts its supplied proofs into
\(A,\neg A\Rightarrow\bot\), obtained by \(\neg L\) from
\(A\Rightarrow\bot,A\). Falsum elimination cuts with the
initial \(\bot\Rightarrow B\).

For RAA, cut the translated premise
\(\Gamma,\neg A\Rightarrow\bot\) with \(\bot\Rightarrow\),
then apply \(\neg R\) to get \(\Gamma\Rightarrow\neg\neg A\).
The Cut/Mix-free kernel \(\neg\neg A\Rightarrow A\) comes from
Id \(A\Rightarrow A\), followed by \(\neg R\) to obtain
\(\Rightarrow A,\neg A\), then \(\neg L\) on its right
\(\neg A\). Cut this kernel into the supplied conclusion.

Universal elimination cuts with
\(\forall xA\Rightarrow A[t/x]\), the instance of \(\forall L\)
over identity of \(A[t/x]\). Universal introduction is \(\forall R\)
with eigenvariable equal to the source's generalized variable \(x\):
the total operation \(A[x/x]\) is literally \(A\), by its recursion,
and \(x\) is free in neither the surviving assumptions nor the
quantified conclusion. Existential introduction is \(\exists R\).
For existential elimination apply \(\exists L\) to the translated
minor premise and cut in the translated major premise. The source
proviso that \(y\) is absent from the surviving assumptions, conclusion
and existential formula implies every registered free-variable proviso.
No additional substitution admissibility theorem is presumed.
ND alpha-conversion cuts with the literal alpha kernel of
Annex~\ref{annex:p7.eqlk.alpha}, rather than treating formula strings
as alpha-classes.

For ND reflexivity use the equality initial and weaken if needed.
For ND replacement cut the two translated premises into the
right-directed kernel \(r=t,A[r/x]\Rightarrow A[t/x]\) from
Annex~\ref{annex:p7.eqlk.replacement}; contract any shared assumptions.
This explicitly supplies the equality translation omitted in the inherited
sequent-completeness sketch and exhausts every ND rule.

The reverse induction returns
\(\operatorname{supp}(\Gamma)\Rightarrow D_\Delta\).
Derive \(D_\Delta\Rightarrow\Delta\) by its finite disjunction tree:
each leaf is contexted identity, each internal node is \(\lor L\),
and the empty-support case is \(\bot L\). Right weakening restores
duplicated succedent occurrences. Cut these two proofs and restore
antecedent multiplicities by left weakening. The exact requested
sequent results, including either empty context.
\end{proof}
\end{annexitem}

\begin{annexitem}{annex:p7.eqlk.completeness}{thm:p07-eqlk-completeness}{Finite-signature reduction and the actual completeness dependency}
\begin{proof}
Soundness was proved in Annex~\ref{annex:p7.eqlk.soundness}.
Suppose \(\Gamma\Rightarrow\Delta\) is valid in the full signature.
Let the smaller signature contain exactly the nonlogical symbols occurring
in this finite sequent. Each formula and term is finite, so this is a
finite signature, including the case with no nonlogical symbols.
Every structure for it expands to a structure for the original signature:
choose one element of its nonempty domain, interpret every unused function
symbol by the constant function with that value (also for nullary symbols),
and every unused relation symbol by the empty relation (also at arity zero).
This is one explicit simultaneous definition using one chosen element,
not a choice from an indexed family of sets. The reduct-invariance lemma
\ref{lem:reduct-invariance} keeps the truth of every formula in the
sequent unchanged. Consequently the sequent is valid already over all
structures and assignments of that finite signature.

Its semantic condition is precisely
\(\operatorname{supp}(\Gamma)\models D_\Delta\), including
\(D_\varnothing=\bot\). Apply the proved identity natural-deduction
completeness result for countable signatures,
\cref{cor:fol-completeness-countable}. This is the actual Henkin theorem
with equality congruence and the quotient term model, not the equality-free
specialization. It produces a finite ND proof of \(D_\Delta\) with
open assumptions in the finite support. The proof is also a proof in
the larger signature: its terms, rules, alpha conditions and eigenvariable
freshness are unchanged. Annex~\ref{annex:p7.eqlk.nd} translates it
to the literal equality sheet and restores the requested context
multiplicities. That construction uses the explicit equality symmetry
and replacement Cuts. Thus it proves the stated cut-admitting
completeness, with no appeal to Cut elimination or to its conclusion.
\end{proof}
\end{annexitem}

\begin{annexitem}{annex:p7.eqlk.cut-failure}{thm:p07-eqlk-cut-failure}{A single essential Cut and the complete backward branch invariant}
\begin{proof}
Let \(a,b\) be distinct constant symbols and choose a variable \(z\)
absent from them. Id gives \(b=a\Rightarrow b=a\).
The primitive \(=L\), with equality \(a=b\) and auxiliary formula
\(b=z\), gives \(a=b,b=b\Rightarrow b=a\). Cut the reflexivity
initial \(\Rightarrow b=b\) against that \(b=b\) occurrence. The
conclusion is \(a=b\Rightarrow b=a\). The two premises of this
single Cut contain neither Cut nor Mix. The conclusion is valid because
actual identity is symmetric.

Now consider any sequent whose antecedent consists only of copies of
\(a=b\) and \(a=a\), and whose succedent consists only of copies
of \(b=a\). Either context may be empty. We prove that every premise
of any applicable Cut/Mix-free inference with such a conclusion has
the same property. This permits a backward induction from the proposed
root through its entire finite derivation tree.

Backward weakening deletes an occurrence; backward contraction duplicates
one. Both preserve the property. No logical connective or quantifier
rule can have such a conclusion, because its principal conclusion
formula has a connective or quantifier while every allowed formula is
an atom. Exchange is multiset identity.

It remains to check \(=L\), including every possible formula and
hole-variable choice. Its displayed equality \(r=t\) is an antecedent
occurrence of the conclusion, so it is literally \(a=a\) or \(a=b\).
In the first case \(r=t=a=a\); the substituted terms are identical,
so the premise formula is exactly the same as the conclusion formula.
In the second case \(r=a,t=b\). The replaced conclusion formula
\(A[b/x]\) is another permitted antecedent atom, either \(a=a\)
or \(a=b\). Substitution preserves its outer constructor, so \(A\)
must itself be an identity atom. Since \(b\) is a distinct constant,
substituting it for a free variable cannot create the constant \(a\)
or remove a function application or any other constant. If
\(A[b/x]=a=a\), neither term can contain the hole, hence
\(A[a/x]=a=a\). If \(A[b/x]=a=b\), the right term is either
the constant \(b\) or the hole variable \(x\), and the left term is
the constant \(a\). Thus \(A[a/x]\) is respectively \(a=b\)
or \(a=a\). The equality rule leaves the succedent unchanged and
removes only its displayed equality occurrence when read backward.
Therefore its premise also has the property. This reasoning excludes
compound formulas, hidden function terms, multiple hole occurrences
and arbitrary hole variables by the literal substitution shape, not
by a subformula theorem for a different calculus.

No sequent with this property is an initial. Contexted Id requires an
identical formula on both sides, but \(b=a\) is syntactically distinct
from \(a=b\) and \(a=a\). Falsum is absent from the antecedent.
The equality initial has a reflexive succedent \(t=t\), which cannot
be \(b=a\) and cannot be empty. The backward invariant hence makes
every possible leaf of a proposed finite Cut/Mix-free proof impossible.
The root \(a=b\Rightarrow b=a\) has the property, so it has no such proof.

For a signature with no two constants, choose two distinct variables
\(u,v\) and use the same three proof steps with a fresh auxiliary
variable. The invariant's allowed atoms become \(u=u,u=v\) on the
left and \(v=u\) on the right. If \(r=t=u=u\), the two substituted
terms remain identical. If \(r=t=u=v\), replacing a hole by \(v\)
can only keep an allowed term or change a displayed \(u=v\) premise
back to \(u=u\): a left output term \(u\) cannot arise from the
hole mapped to the distinct \(v\), while a right output \(v\) is
either retained \(v\) or that hole. An output \(u=u\) has no hole
occurrence. The same initial exclusions follow. This checks the
variable case even when the hole variable equals \(u\) or \(v\),
since the total operation is a single substitution, not recursive
replacement of the inserted variable.

This is directly a failure of Cut admissibility: its two free premises
exist and its free conclusion does not. Mix with one selected occurrence
on each side is the same inference, so it fails admissibility as well.
Finally the root is valid but has no free proof, establishing the
failure of semantic completeness of the fragment.
\end{proof}
\end{annexitem}
